\chapter{La reproduction humaine}

\noindent\textit{La reproduction humaine repose sur le fonctionnement coordonné de deux
appareils reproducteurs, masculin et féminin, dont l'activité est déclenchée à la
puberté. Ce chapitre en présente l'anatomie, le cycle féminin, la fécondation et le
déroulement de la grossesse.}
\vspace{8pt}

\connecteBanner

\section{Les appareils reproducteurs}

\begin{propriete}[Appareil reproducteur masculin]
Les \textbf{testicules} produisent les \textbf{spermatozoïdes} ainsi qu'une hormone, la
\textbf{testostérone}. Les spermatozoïdes transitent par l'\textbf{épididyme} puis le
\textbf{canal déférent}, se mélangent aux sécrétions des glandes annexes (vésicules
séminales, prostate) pour former le \textbf{sperme}, expulsé par l'urètre lors de
l'éjaculation.
\end{propriete}

\begin{propriete}[Appareil reproducteur féminin]
Les \textbf{ovaires} produisent les \textbf{ovules} ainsi que des hormones (œstrogènes,
progestérone). Un ovule libéré rejoint une \textbf{trompe utérine}, qui le conduit vers
l'\textbf{utérus}, organe dans lequel se développe un éventuel embryon. Le \textbf{vagin}
relie l'utérus à l'extérieur du corps.
\end{propriete}

\begin{remarque}
C'est à la \textbf{puberté} que ces deux appareils deviennent fonctionnels, accompagnés
de l'apparition des caractères sexuels secondaires (pilosité, mue de la voix, développement
des seins\dots).
\end{remarque}

\section{Le cycle féminin}

\begin{definition}[Cycle menstruel]
Le \textbf{cycle menstruel} dure en moyenne \textbf{28 jours}. Il débute le premier jour
des \textbf{règles} (menstruations : évacuation de la paroi utérine en l'absence de
grossesse) et se termine la veille des règles suivantes.
\end{definition}

\begin{propriete}[L'ovulation]
Vers le \textbf{14\textsuperscript{e} jour} d'un cycle de 28 jours (soit environ au
milieu du cycle), un ovaire libère un ovule : c'est l'\textbf{ovulation}. C'est la période
du cycle où une fécondation est possible.
\end{propriete}

\begin{illus}
\begin{tikzpicture}[scale=0.8]
\draw[->,onbmuted,line width=1.2pt] (0,0) -- (11.5,0);
\filldraw[onbo,line width=0pt] (0.3,0) circle(0.09);
\node[below,font=\scriptsize] at (0.3,-0.25) {jour 1 : règles};
\filldraw[onbgreen,line width=0pt] (5.5,0) circle(0.09);
\node[below,font=\scriptsize] at (5.5,-0.25) {jour 14 : ovulation};
\filldraw[onbo2,line width=0pt] (10.7,0) circle(0.09);
\node[below,font=\scriptsize] at (10.7,-0.25) {jour 28 : fin du cycle};
\end{tikzpicture}
\fig{Figure — Le cycle menstruel : de 28 jours en moyenne, ovulation vers le jour 14.}
\end{illus}

\begin{exemple}
Pour un cycle de $30$ jours, l'ovulation a lieu vers le jour $30\div2=15$.
\end{exemple}

\section{La fécondation}

\begin{definition}[Fécondation]
La \textbf{fécondation} est la fusion d'un \textbf{spermatozoïde} et d'un \textbf{ovule},
qui a lieu dans une \textbf{trompe utérine}. Elle donne naissance à une nouvelle cellule,
la \textbf{cellule-œuf}, qui contient l'information génétique des deux parents.
\end{definition}

\begin{illus}
\begin{tikzpicture}[scale=0.8]
\filldraw[onbo!20,draw=onbo,line width=1.1pt] (0,0) rectangle (5.2,0.9);
\node[font=\small] at (2.6,0.45) {spermatozoïde $+$ ovule};
\draw[->,onbgreen,line width=1.3pt] (5.2,0.45) -- (6.6,0.45);
\node[above,onbgreen,font=\scriptsize] at (5.9,0.65) {fécondation};
\filldraw[onbgreen!20,draw=onbgreen,line width=1.1pt] (6.8,0) rectangle (9.6,0.9);
\node[font=\small] at (8.2,0.45) {cellule-œuf};
\end{tikzpicture}
\fig{Figure — La fécondation : fusion d'un spermatozoïde et d'un ovule.}
\end{illus}

\section{La grossesse}

\begin{propriete}[De la fécondation à la naissance]
Après la fécondation, la cellule-œuf se divise plusieurs fois en migrant vers l'utérus,
où elle s'implante dans la paroi utérine : c'est la \textbf{nidation}. L'embryon, puis le
\textbf{fœtus}, se développe ensuite pendant une \textbf{grossesse} d'environ \textbf{9
mois} (soit environ $39$ semaines).
\end{propriete}

\begin{definition}[Placenta]
Le \textbf{placenta} est un organe qui se forme au contact de la paroi utérine. Il permet
les \textbf{échanges} entre le sang de la mère et celui du fœtus (nutriments, dioxygène,
déchets), \textbf{sans que les deux sangs ne se mélangent}.
\end{definition}

\section{La contraception}

\begin{propriete}[Éviter une grossesse non désirée]
Différentes méthodes de \textbf{contraception} permettent d'éviter une grossesse non
désirée : le \textbf{préservatif} (qui empêche en plus la transmission de certaines
infections), la \textbf{pilule contraceptive} (qui bloque l'ovulation), et d'autres
méthodes hormonales ou mécaniques.
\end{propriete}

% ═══════════════════════════════════════════════════════════════
\section{L'essentiel à retenir}

\begin{synthese}
\textbf{Appareil masculin.} Testicules (spermatozoïdes, testostérone) $\to$ épididyme
$\to$ canal déférent.\\[4pt]
\textbf{Appareil féminin.} Ovaires (ovules, hormones) $\to$ trompes $\to$ utérus.\\[4pt]
\textbf{Cycle menstruel.} $28$ jours en moyenne ; ovulation vers le jour $14$.\\[4pt]
\textbf{Fécondation.} Spermatozoïde $+$ ovule $\to$ cellule-œuf, dans une trompe
utérine.\\[4pt]
\textbf{Grossesse.} Nidation dans l'utérus ; développement d'environ $9$ mois ; échanges
mère-fœtus via le placenta.\\[4pt]
\textbf{Piège fréquent.} Le sang de la mère et celui du fœtus ne se mélangent
\textbf{jamais} : les échanges se font à travers la paroi du placenta.
\end{synthese}

% ═══════════════════════════════════════════════════════════════
\section{Méthodes \& savoir-faire}

\methode{Retracer le trajet des spermatozoïdes ou de l'ovule}
Spermatozoïdes : testicules $\to$ épididyme $\to$ canal déférent. Ovule : ovaire $\to$
trompe $\to$ utérus.
\begin{exemple}
Un spermatozoïde ne peut pas atteindre l'ovule sans traverser l'épididyme puis le canal
déférent.
\end{exemple}

\methode{Calculer le jour approximatif de l'ovulation}
Diviser la durée du cycle par deux.
\begin{exemple}
Cycle de $26$ jours : ovulation vers le jour $26\div2=13$.
\end{exemple}

\methode{Distinguer fécondation et nidation}
La fécondation a lieu dans la \textbf{trompe} (fusion des gamètes) ; la nidation a lieu
dans l'\textbf{utérus} (implantation de l'embryon), quelques jours plus tard.
\begin{exemple}
La cellule-œuf se forme dans la trompe avant de migrer vers l'utérus pour s'y implanter.
\end{exemple}

% ═══════════════════════════════════════════════════════════════
\section{Exercices d'application}
\begin{multicols}{2}

\rubrique{Appareils reproducteurs}
\exo{}\diff{1} Quel organe produit les spermatozoïdes ?

\exo{}\diff{1} Quel organe produit les ovules ?

\exo{}\diff{2} Quel est le rôle de l'utérus ?

\exo{}\diff{2} À quel moment de la vie les appareils reproducteurs deviennent-ils
fonctionnels ?

\rubrique{Cycle féminin}
\exo{}\diff{1} Quelle est la durée moyenne d'un cycle menstruel ?

\exo{}\diff{2} Qu'est-ce que l'ovulation ?

\exo{}\diff{2} Pour un cycle de $30$ jours, calculer le jour approximatif de
l'ovulation.

\rubrique{Fécondation et grossesse}
\exo{}\diff{1} Où a lieu la fécondation ?

\exo{}\diff{2} Qu'est-ce que la nidation ?

\exo{}\diff{2} Quel est le rôle du placenta ?

\exo{}\diff{1} Quelle est la durée moyenne d'une grossesse ?

% ═══════════════════════════════════════════════════════════════
\end{multicols}

\section{Exercices d'entraînement}
\begin{multicols}{2}

\rubrique{Appareils reproducteurs}
\exo{}\diff{2} Citer, dans l'ordre, les organes traversés par les spermatozoïdes avant
l'éjaculation.

\exo{}\diff{3} Quelles hormones sont produites par les ovaires ?

\rubrique{Cycle féminin}
\exo{}\diff{2} Pour un cycle de $26$ jours, calculer le jour approximatif de
l'ovulation.

\exo{}\diff{3} Pourquoi les règles surviennent-elles en l'absence de fécondation ?

\exo{}\diff{3} Une femme a un cycle de $28$ jours et ses dernières règles ont commencé
le $1$\textsuperscript{er} jour du mois. Vers quel jour du mois aura probablement lieu son
ovulation ?

\rubrique{Fécondation et grossesse}
\exo{}\diff{3} Expliquer la différence entre fécondation et nidation.

\exo{}\diff{3} Pourquoi le sang maternel et le sang fœtal ne se mélangent-ils jamais au
niveau du placenta ?

\exo{}\diff{2} Une grossesse dure environ $39$ semaines. Convertir cette durée en jours,
puis en mois (arrondir).

\rubrique{Contraception}
\exo{}\diff{2} Citer deux méthodes de contraception.

\exo{}\diff{3} Quel est l'avantage particulier du préservatif par rapport à la pilule
contraceptive ?

% ═══════════════════════════════════════════════════════════════
\end{multicols}

\section{Sujets type BEPC}

\sujetbac[Les appareils reproducteurs]
\begin{enumerate}[label=\arabic*)]
\item Cite l'organe qui produit les spermatozoïdes et celui qui produit les ovules.
\item Cite, dans l'ordre, les organes traversés par un ovule libéré jusqu'à l'utérus.
\item Quel est le rôle de l'utérus ?
\item À quel moment de la vie ces appareils deviennent-ils fonctionnels ?
\end{enumerate}

\sujetbac[Cycle menstruel et fécondation]
\begin{enumerate}[label=\arabic*)]
\item Donne la durée moyenne d'un cycle menstruel.
\item Pour un cycle de $28$ jours, calcule le jour approximatif de l'ovulation.
\item Où a lieu la fécondation ?
\item Que produit la fécondation ?
\end{enumerate}

\sujetbac[La grossesse]
\begin{enumerate}[label=\arabic*)]
\item Qu'est-ce que la nidation, et où a-t-elle lieu ?
\item Quelle est la durée moyenne d'une grossesse humaine, en mois puis en semaines ?
\item Quel est le rôle du placenta ?
\item Pourquoi le placenta permet-il des échanges sans mélanger les deux sangs ?
\end{enumerate}

% ═══════════════════════════════════════════════════════════════
\section{Corrigés}
\begin{multicols}{2}

\corrige{de l'exercice 1}
Les testicules.

\corrige{de l'exercice 2}
Les ovaires.

\corrige{de l'exercice 3}
Il accueille le développement de l'embryon puis du fœtus lors d'une grossesse.

\corrige{de l'exercice 4}
À la puberté.

\corrige{de l'exercice 5}
$28$ jours en moyenne.

\corrige{de l'exercice 6}
La libération d'un ovule par un ovaire.

\corrige{de l'exercice 7}
$30\div2=15$ : vers le jour $15$.

\corrige{de l'exercice 8}
Dans une trompe utérine.

\corrige{de l'exercice 9}
L'implantation de la cellule-œuf (en division) dans la paroi de l'utérus.

\corrige{de l'exercice 10}
Il permet les échanges entre le sang maternel et le sang fœtal (nutriments, dioxygène,
déchets), sans mélange des deux sangs.

\corrige{de l'exercice 11}
Environ $9$ mois.

\corrige{de l'exercice 12}
Testicules, épididyme, canal déférent.

\corrige{de l'exercice 13}
Les œstrogènes et la progestérone.

\corrige{de l'exercice 14}
$26\div2=13$ : vers le jour $13$.

\corrige{de l'exercice 15}
Car la paroi utérine, préparée pour accueillir un éventuel embryon, est évacuée en
l'absence de fécondation.

\corrige{de l'exercice 16}
Vers le $14$\textsuperscript{e} jour du mois (jour $1+14-1=14$).

\corrige{de l'exercice 17}
La fécondation est la fusion du spermatozoïde et de l'ovule dans la trompe ; la nidation
est l'implantation, quelques jours plus tard, de la cellule-œuf (en division) dans la
paroi de l'utérus.

\corrige{de l'exercice 18}
Car les échanges se font à travers la paroi du placenta, qui laisse passer certaines
substances sans mettre directement en contact les deux circulations sanguines.

\corrige{de l'exercice 19}
$39\times7=273$ jours, soit environ $273\div30\approx9$ mois.

\corrige{de l'exercice 20}
Par exemple : le préservatif et la pilule contraceptive.

\corrige{de l'exercice 21}
Il protège en plus contre la transmission de certaines infections sexuellement
transmissibles.

\corrige{du sujet 1}
1) Les testicules produisent les spermatozoïdes ; les ovaires produisent les ovules.\\
2) Ovaire, trompe utérine, utérus.\\
3) Il accueille le développement de l'embryon puis du fœtus.\\
4) À la puberté.

\corrige{du sujet 2}
1) $28$ jours en moyenne.\\
2) $28\div2=14$ : vers le jour $14$.\\
3) Dans une trompe utérine.\\
4) Une cellule-œuf.

\corrige{du sujet 3}
1) L'implantation de la cellule-œuf dans la paroi de l'utérus.\\
2) Environ $9$ mois, soit environ $39$ semaines.\\
3) Il permet les échanges entre le sang maternel et le sang fœtal (nutriments,
dioxygène, déchets).\\
4) Car les échanges se font à travers la paroi du placenta, qui agit comme une barrière
sélective entre les deux circulations sanguines.

\end{multicols}
\exercicesBanner
